\section{A Glimpse at the Data}



\textbf{Analysis of Data Formats:}

The data pipeline begins with two types of raw text files for each simulation instance: one specifying the initial conditions and another containing the time-series output from the Brutus N-body integrator.

\paragraph{Initial Conditions (\url{ic_hierarchical_*.txt})}
This is a 3-line text file containing 21 space-delimited, double-precision floating-point numbers. 
These numbers represent the complete initial state vector of the three-body system at time $t=0$. 
The structure is a flattened array corresponding to the seven state variables for each of the three bodies for each line. Together,
it can be represented as: 
\texttt{[m1, x1, y1, z1, vx1, vy1, vz1, m2, x2, y2, z2, vx2, vy2, vz2, m3, x3, y3, z3, vx3, vy3, vz3]}. 
This file is the primary source for calculating the conditional input vector for the PINN.

\begin{lstlisting}[
    language=python,
    caption={Sample Initial Condition File (\texttt{ic\_hierarchical\_0000.txt})},
    label={lst:ic_sample_1b},
    frame=single,
    backgroundcolor=\color{lightgray},
    basicstyle=\ttfamily\small,
    breaklines=true
]
1.3490330328149076 -0.9689510833073296 0.007710980988132066 -0.015155493143494857 0.003082579960339502 -0.9723812422861391 0.01317780514703599 0.3068019185652623 -0.6619848581497906 -0.005365500630046534 0.0070754473066994734 0.002840479118371817 1.3293141254107488 -0.013692746386640507 1.3441650486198304 1.1354206743022845 -0.0016313630906544987 0.002882937931195984 0.00441920746565285 0.6549777259648452 0.001846566609351255
\end{lstlisting}

\paragraph{Brutus Simulation Output (\url{brutus_hierarchical_*.txt})}
This is a time-series data file that provides the ground-truth 
trajectories for the data-fidelity loss term, $\mathcal{L}_{data}$. The file is structured as a sequence of blocks, 
where each block corresponds to a specific snapshot in time. Each block begins with a header line containing three values: 
the simulation time \texttt{t}, the number of bodies \texttt{N} (which is always 3), and \texttt{t\_cpu} (CPU Time). 
This header is followed by \texttt{N} lines, where each line contains the seven state variables 
(\texttt{m, x, y, z, vx, vy, vz}) for one of the three bodies at that specific time \texttt{t}.

Thus, every Brutus output file has a repetitive block format that contains:
\begin{itemize}
    \item Header: \texttt{t, N, t\_cpu}
    \item Body 1: \texttt{m1, x1, y1, z1, vx1, vy1, vz1}
    \item Body 2: \texttt{m2, x2, y2, z2, vx2, vy2, vz2}
    \item Body 3: \texttt{m3, x3, y3, z3, vx3, vy3, vz3}
\end{itemize}

A raw sample output from a Brutus simulation file is shown below, illustrating the data format for a few time steps:

\begin{lstlisting}[
    caption={Sample Raw Output from Brutus Simulation File},
    label={lst:brutus_output},
    frame=single,
    backgroundcolor=\color{lightgray},
    basicstyle=\ttfamily\small,
    breaklines=true,
    columns=flexible,
    xleftmargin=0.5cm,
    xrightmargin=0.5cm,
    aboveskip=0.5cm,
    belowskip=0.5cm,
    numbers=left,
    numberstyle=\tiny\color{gray}
]
0 3 0
1.3490330328149076 -0.9689510833073296 0.0077109809881320658 -0.015155493143494857 0.0030825799603395022 -0.97238124228613909 0.01317780514703599
0.30680191856526229 -0.66198485814979058 -0.0053655006300465337 0.0070754473066994734 0.002840479118371817 1.3293141254107488 -0.013692746386640507
1.3441650486198304 1.1354206743022845 -0.0016313630906544987 0.002882937931195984 0.0044192074656528501 0.65497772596484516 0.001846566609351255

0.080000000000000002 3 2.17338
1.3490330328149076 -0.95746912715152098388870457096922081377499 -0.068395326027771184330918692421700212149959 -0.013374646173250254126208350188876903106782 0.27927408182448689800114515637261157892989 -0.90345482920569108996083189130475635165807 0.030685786229670570110204639660396401512047
0.30680191856526229 -0.70556381350356795672897465039496696390801 0.093633356547250511770656254567903108831681 0.002818445191678544590613756190767057248414 -1.0718184960048500457752700625063898462062 1.0286823897602963737557930793641611624554 -0.089839191865848882805429923629505460110143
1.3441650486198304 1.1344968129380317254090080254393242791445 0.050753329467245612450222255049915160172 0.0030230315435012235246815871672795058650157 -0.027484630586336917599520346861990470320003 0.65442004534327653581772256816764955369181 0.0016553927250949034582096866402066475866725
\end{lstlisting}





In the following sections, we lay out a detailed, three-phase implementation plan for the CTL-PINN training methodology.  
This guide is intended to be sufficiently complete for reproducibility and practical deployment of the model.  
We assume the datasets (hierarchical and quasi-random simulations) have been generated as described in Section 2 and the supporting data pipeline document.  
The implementation is broken into: Phase 1: Data preparation and curriculum setup, Phase 2: The CTL-PINN training loop, and Phase 3: Active learning-based data enhancement. 
